\section{Metodología experimental}
%=====================================================================

\textbf{Datos.} Se emplearon los repositorios públicos de PhysioNet
\cite{goldberger2000}: la \emph{MIT-BIH Arrhythmia Database}
\cite{moody2001} para la clasificación de latidos (4 clases AAMI: N, S,
V, F), y las bases PTB Diagnostic ECG (PTBDB) y PTB-XL \cite{wagner2020}
para validar la lectura de registros de 12 derivaciones con metadatos
heterogéneos. La partición de evaluación sigue la propuesta de de Chazal
et al.~\cite{dechazal2004} (DS1/DS2 por registro). Los registros de
MIT-BIH se leyeron directamente de los archivos \texttt{.hea}/\texttt{.dat}/\texttt{.atr}
con el lector WFDB descrito en la Sección~2; los de PTBDB y PTB-XL se
usaron para comprobar la conversión de unidades, la ganancia por canal y
el número variable de derivaciones. Adicionalmente, el generador
sintético incluido en el proyecto permitió verificar la cadena completa
(cabecera, formato 16/212, \emph{checksum} y anotaciones con SKIP) en
casos controlados.

\textbf{Experimento A (detección).} Se evaluó el detector Pan-Tompkins
por derivación y la fusión por consenso sobre registros reales de
MIT-BIH (360\,Hz, 2 derivaciones, formato 212) con tolerancia
$\tau=0{,}15$\,s, comparando contra las anotaciones de referencia del
\texttt{.atr}. La validación se realizó registro a registro.

%\textbf{Experimento B (escalabilidad).} Registros de 12 derivaciones a
%500\,Hz con duración de 30, 60, 120, 300 y 600\,s, en formatos 16 y 212.
%Para cada uno se midieron la carga completa (cabecera, validación,
%decodificación, conversión y anotaciones), la detección en las 12
%derivaciones con fusión y la generación de la impresión estándar de
%5\,s. Se realizaron dos repeticiones de calentamiento y cinco medidas,
%con recolección de basura previa a cada una; se informa la media.

\textbf{Campaña de 15 pruebas preliminares (esta entrega).} Con el
clasificador ya entrenado con DS1, se ejecutó el flujo completo
(carga $\rightarrow$ detección QRS $\rightarrow$ fusión $\rightarrow$
extracción de 257 características $\rightarrow$ clasificación con el
bosque $\rightarrow$ validación contra referencia) sobre 15 registros
de MIT-BIH: 200, 201, 202, 203, 205, 207, 208, 209, 210, 212, 213, 214,
215, 217 y 219. Cada prueba corresponde a un registro (P001--P015 en la
Tabla~\ref{tab:campana}) y se ejecutó en una sesión independiente de la
JVM. Para cada prueba se registraron: el número de latidos por clase en
la referencia (Soporte N/S/V/F), la matriz de confusión $4\times4$, las
métricas por clase (precisión, \emph{recall}, especificidad, $F_1$) y
las métricas globales (\emph{accuracy}, $F_1$ macro, \emph{balanced
accuracy}), además de los tiempos y el consumo de recursos de cada
operación del flujo (carga, decodificación, análisis QRS, clasificación
RF, apertura del visor y exportación PNG). La Tabla~\ref{tab:campana}
resume los resultados de clasificación; la Tabla~\ref{tab:tiempos}
resume los tiempos y el heap máximo por operación, promediados sobre las
15 pruebas. Los informes PDF individuales
(\texttt{informe\_rendimiento\_*.pdf}) se generaron con el módulo de
rendimiento descrito en la Sección~\ref{sec:rend}.

%\textbf{Entorno.} Las pruebas de clasificación y de rendimiento de esta
%entrega se ejecutaron en Windows~11 (amd64) con 20 procesadores,
%OpenJDK~25.0.4.1 (Eclipse Adoptium) y \texttt{-Xmx2g}. El experimento B
%de escalabilidad se ejecutó en Linux~6.18 (amd64), OpenJDK~21.0.12, un
%único procesador y 4\,GB de RAM, en modo sin pantalla; por tanto, sus
%tiempos no son directamente comparables con los de la
%Tabla~\ref{tab:tiempos}. La verificación de la generación de PDF se
%realizó con PDFBox~2.0.29; el código está escrito para la serie 3.x
%declarada en el \texttt{pom.xml} (solo difiere la creación de las
%fuentes).

\textbf{Pruebas automáticas.} 31 pruebas JUnit~5 cubren la lectura
(cabeceras, formatos 16/212/8/24/32 con vectores construidos a mano,
conversión de unidades, anotaciones con SKIP, \emph{checksum}, ausencia
de ganancia), el procesamiento (detección, fusión, validación, RR/FC),
la visualización (distribuciones, zoom con punto fijo, selección y
cursor), la exportación (calibración exacta, coincidencia píxel a píxel
con la pantalla, paginación) y el rendimiento (medidores, anidamiento,
PDF y evaluación experimental).
